\chapter{Prijateljske funkcije i klase te operatori}

\section{Prijateljske funkcije i klase}

Sadržaj \textit{Klijent.h} datoteke:

\begin{minted}{C++}
#pragma once
#include <string>
#include <vector>

class Klijent {
private:
    std::string ime;
    std::vector<double> transakcije;
    double ukupno();
public:
    Klijent(std::string i = "nepoznat", double pocetno = 0);
    void uplati(double iznos, Klijent &k);
    void info();
};
\end{minted}


Sadržaj \textit{Klijent.cpp} datoteke:

\begin{minted}{C++}
#include <iostream>
#include <string>
#include "Klijent.h"

using namespace std;

Klijent::Klijent(string i, double p)
	: ime(i), transakcije(1,p) {}

void Klijent::uplati(double iznos, Klijent &k) {
    k.transakcije.push_back(iznos);
    transakcije.push_back(-iznos);
}
	
double Klijent::ukupno() {
	double zbroj = 0;
	vector<double>::iterator it;
	for(it = transakcije.begin();
		it != transakcije.end(); ++it)
			zbroj += *it;
	return zbroj;
}

void Klijent::info() {
	cout << "Podaci o racunu klijenta:" << endl;
	cout << "\tIme: " << ime << endl;
	cout << "\tTransakcije: " << endl;
	if(transakcije.empty())
		cout << "\tNema transakcija!" << endl;
	else {
		vector<double>::iterator it;
		for(it = transakcije.begin();
			it != transakcije.end(); ++it)
				cout << "\t\t" << *it 
					 << " EUR" << endl;
	}
	cout << "Ukupno na racunu: "
		 << ukupno() << " EUR." << endl;
}
\end{minted}


Sadržaj \textit{main.cpp} datoteke:

\begin{minted}{C++}
#include <iostream>
#include "Klijent.h"

using namespace std;

int main(void) {
	Klijent a, b("Bob",1200);
	a.uplati(1150.99,b);
	a.info();
	b.info();
	
	return 0;
}
\end{minted}


\textbf{Prijateljske funkcije} (engl.\ \textit{friend functions}) 
predstavljaju posebne funkcije koje nisu članice klase, ali im je dopušten pristup privatnim i zaštićenim članovima te klase. 
Time se zadržava enkapsulacija, uz kontrolirano dopuštanje pristupa podacima kada je to potrebno.
Takve funkcije deklariraju se unutar definicije klase pomoću ključne riječi \texttt{friend}.
\medskip

Važno je naglasiti da \texttt{friend} funkcija nije metoda klase.
Primjerice,
ako imamo (globalnu) funkciju prototipa \texttt{int funkcija(int broj);}
koja koristi neki privatan član neke klase \texttt{Klasa},
u definiciji te klase deklariramo tu funkciju kao prijateljsku ovako:
\begin{minted}{C++}
class Klasa {
friend int funkcija(int broj);
...
};
\end{minted}
Međutim, ako je objekt \texttt{obj} tipa klase \texttt{Klasa},
ne možemo pisati
\begin{minted}{C++}
int rezultat = obj.funkcija(123);
\end{minted}
Pri pokušaju kompajliranja prethodne linije koda,
kompajler će javiti grešku da klasa \texttt{Klijent} nema člana s imenom \texttt{provjera\_uvjeta}.
\medskip

Deklaracija prijateljske funkcije može se nalaziti u bilo kojem dijelu klase, 
obzirom da na njih ne utječe kontrola pristupa
(tj.\ modifikatori \texttt{public}, \texttt{private} i \texttt{protected}).

\begin{zadatak}\label{zad:klijent-friend}
    U \textcolor{red}{zadatku ...} implementirana je klasa \texttt{Klijent}
    pomoću koje banka pamti transakcije pojedinog klijenta u vektoru realnih brojeva.
    Nenegativne iznose, tj.\ realne brojeve, u tom vektoru nazivamo primanjima tog klijenta.
    Napišite (globalnu) funkciju 
    \begin{center}
        \texttt{bool provjera\_uvjeta(double rata, Klijent\& k)}
    \end{center}
    koja vraća \texttt{true} ako i samo ako klijent \texttt{k} 
    može podignuti u toj banci kredit s ratom koja iznosi \texttt{rata} eura.
    Klijent može podignuti s tim iznosom rate ako
    mu iznos te rate iznosi najviše 45\% prosječnog primanja.
    \medskip

    Napišite i program za testiranje u \textit{main.cpp} datoteci.
    Kako bi ta datoteka bila pregledna, funkciju \texttt{provjera\_uvjeta}
    implementirajte u \textit{Klijent.cpp} datoteci.
    Zbog kasnije ponovne upotrebe koda, prvo klasi \texttt{Klijent}
    dodajte pomoćnu privatnu metodu \texttt{double prosjek\_primanja();}
    koja vraća iznos prosječnog primanja klijenta.
\end{zadatak}


\begin{rjesenje}
Dodajemo klasi \texttt{Klijent}
(u \textit{Klijent.h} datoteci) deklaraciju tražene pomoćne privatne metode \texttt{prosjek\_primanja}:

\begin{minted}[
        linenos,
        numbersep=5pt
    ]{C++}
class Klijent {
private:
    double prosjek_primanja();
    ...
}
\end{minted}

Dodajemo implementaciju te metode u \textit{Klijent.cpp} datoteku:

\begin{minted}[
        linenos,
        numbersep=5pt
    ]{C++}
double Klijent::prosjek_primanja() {
    double prosjek = 0;
    int brojac = 0;
    vector<double>::iterator it;
    for(it = transakcije.begin(); 
        it != transakcije.end(); ++it) {
            if(*it > 0) {
                prosjek += *it;
                brojac++;
            }
    }
    if(brojac != 0)
        prosjek /= brojac;
    return prosjek;
}
\end{minted}

Upotrebom te metode, implementiramo funkciju \texttt{provjera\_uvjeta}
u \texttt{Klijent.cpp} datoteci:

\begin{minted}[
        linenos,
        numbersep=5pt
    ]{C++}
bool provjera_uvjeta(double rata, Klijent& k){
    return 0.45 * k.prosjek_primanja() >= rata;
}
\end{minted}    

Obzirom da je metoda \texttt{prosjek\_primanja}
u klasi \texttt{Klijent} označena kao privatna,
trebamo u klasi \texttt{Klijent} dati dopuštenje 
funkciji \texttt{provjera\_uvjeta}
za upotrebu privatnih dijelova klase.
To činimo označavanjem te funkcije kao prijateljske u klasi \texttt{Klijent}
(u \textit{Klijent.h} datoteci):

\begin{minted}[
        linenos,
        numbersep=5pt
    ]{C++}
class Klijent {
friend bool provjera_uvjeta(double rata, Klijent &k);
...
};
\end{minted}

Sadržaj \textit{main.cpp} datoteke za testiranje:

\begin{minted}[
        linenos,
        numbersep=5pt
    ]{C++}
#include <iostream>
#include "Klijent.h"

using namespace std;

int main(void) {
	Klijent a, b("Bob",1200);
	cout << provjera_uvjeta(535,b) << endl;
	a.uplati(1150.99,b);
	cout << provjera_uvjeta(535,b) << endl;
	a.info();
	b.info();
	
	return 0;
}
\end{minted}
\end{rjesenje}


Osim funkcija, također i cijela klasa može biti deklarirana kao prijatelj druge klase (tj.\ kao \textbf{prijateljska klasa}). 
Time se svim članicama prijateljske klase omogućuje pristup privatnim i zaštićenim članovima klase koja joj je dodijelila status prijatelja.
\medskip

Prijateljska klasa također se deklarira pomoću ključne riječi \texttt{friend} unutar definicije druge klase.
Primjerice, ako želimo nekoj klasi \texttt{Banka}
dopustiti pristup svim članovima klase \texttt{Racun},
u definiciji te klase \texttt{Racun}
navodimo deklaraciju klase \texttt{Banka} pomoću ključne riječi \texttt{friend} ovako:
\begin{minted}{C++}
class Racun {
friend class Banka;
...
};
\end{minted}


Istaknimo da prijateljstvo klasa nije uzajamno, ne nasljeđuje se i nije tranzitivno.
Primjerice, u slučaju gornjih klasa \texttt{Racun} i \texttt{Banka}, iako je klasa \texttt{Banka}
prijateljska klasa klasi \texttt{Racun},
to ne znači da je klasa \texttt{Racun} prijateljska
klasi \texttt{Banka}.
Klasa \texttt{Banka} također nije prijateljska
klasama koje nasljeđuju klasu \texttt{Racun}.
Nadalje,
ako je neka klasa prijateljska klasi \texttt{Banka},
ona nije prijateljska klasa klasi \texttt{Racun}.


\begin{zadatak}\label{zad:kredit}
    Napišite klasu \texttt{Kredit} koja modelira bankovni kredit. 
    Klasa treba sadržavati podatke
    o glavnici (tj.\ iznosu kredita u eurima),
    godišnjoj kamatnoj stopi (izraženoj u postocima)
    i roku otplate (u mjesecima).
    Konstruktor treba primati vrijednosti za 
    ta tri podatka te njima inicijalizirati objekt.
    Dodajte klasi metodu \texttt{mjesecna\_rata}
    koja ne prima dodatne argumente, nego samo izračunava i vraća 
    iznos mjesečne anuitetske rate.
    Ta rata se računa prema formuli
    $$\frac{G\cdot m\cdot(1+m)^n}{(1+m)^n-1},$$
    gdje je $G$ glavnica kredita, $m$ mjesečna kamatna stopa
    te $n$ broj mjeseci otplate.
    \medskip

    Zatim dodajte toj klasi javnu metodu 
    \begin{center}
        \texttt{bool provjera\_sposobnosti(Klijent \&k)}
    \end{center}
    koja za klijenta tipa klase \texttt{Klijent}
    (iz zadatka \ref{zad:klijent-friend})
    provjerava njegovu kreditnu sposobnost,
    tj.\ može li taj klijent dići taj kredit
    (pri čemu je jedini uvjet obzirom na primanja i iznos rate kredita opisan u zadatku \ref{zad:klijent-friend}).
    Navedena metoda mora koristiti metode klase \texttt{Klijent},
    pri čemu joj se to mora omogućiti upotrebom \texttt{friend} deklaracije odgovarajuće klase.
    \medskip

    Naposlijetku, napišite i program za testiranje.
\end{zadatak}


\begin{rjesenje}
    Deklaracija klase \texttt{Kredit} (u datoteci \textit{Kredit.h}):
    \begin{minted}[
        linenos,
        numbersep=5pt
    ]{C++}
#pragma once
#include "Klijent.h"

class Kredit {
	double glavnica;
	double stopa;
	unsigned rok;
public:
    Kredit(double g, double s, unsigned r);
    double mjesecna_rata();
    bool provjera_sposobnosti(Klijent& k);	
};
    \end{minted}

Implementacija funkcija članica klase u \textit{Kredit.cpp} datoteci,
pri čemu smo zbog upotrebe funkcije za potenciranje \texttt{pow}
uključili biblioteku matematičkih funkcija \texttt{cmath}:
    \begin{minted}[
        linenos,
        numbersep=5pt
    ]{C++}
#include <cmath>
#include "Klijent.h"
#include "Kredit.h"

using namespace std;

Kredit::Kredit(double g, double s, unsigned r)
    : glavnica(g), stopa(s), rok(r) {}

double Kredit::mjesecna_rata() {
    double m = stopa / 100 / 12;
    return (glavnica * m * pow(1+m,rok))
		/ (pow(1+m,rok)-1);
}

bool Kredit::provjera_sposobnosti(Klijent& k) {
    return k.prosjek_primanja() * 0.45 >= mjesecna_rata();
}
    \end{minted}
    U funkciji \texttt{Kredit::mjesecna\_rata} podijelili
    smo godišnju stopu \texttt{stopa} s brojem $12$
    kako bi dobili iznos mjesečne stope
    te s brojem $100$ kako bismo postotak pretvorili u 
    realan broj između 0 i 1.
    \medskip
    
    Uočimo da metoda \texttt{provjera\_sposobnosti} klase \texttt{Kredit} pristupa metodi \texttt{prosjek\_primanja} klase \texttt{Klijent}
    koja je u toj klasi označena kao privatna.
    Prema tome, potrebno je u definiciji klase \texttt{Klijent}
    (u \textit{Klijent.h} datoteci)
    deklarirati klasu \texttt{Kredit} kao prijateljsku klasu
    (kako bi metode klase \texttt{Kredit} dobile pristup privatnim članovima klase \texttt{Klijent}):
    \begin{minted}[
        linenos,
        numbersep=5pt
    ]{C++}
    class Klijent {
    friend class Kredit;    
    ...
    };
    \end{minted}

    Sadržaj \textit{main.cpp} datoteke s primjerom programa za testiranje:

    \begin{minted}[
        linenos,
        numbersep=5pt
    ]{C++}
#include <iostream>

#include "Klijent.h"
#include "Kredit.h"

using namespace std;

int main(void) {
    Klijent a, b("Bob",1200);
    a.uplati(1150.99,b);
	
    Kredit stan(130000,2.85,360);
    cout << stan.mjesecna_rata() << endl;
    
    return 0;
}
\end{minted}
\end{rjesenje}


Pažljiviji čitatelji uočit će da u datoteci \textit{Klijent.h} u rješenju prethodnog zadatka 
nismo uključili datoteku
\textit{Kredit.h} iako u toj datoteci imamo sljedeći kod:
    \begin{minted}[
        linenos,
        numbersep=5pt
    ]{C++}
class Klijent {
friend class Kredit;    
...
};
    \end{minted}
Naime, nije potrebno uključiti zaglavlje neke klase samo zato što se njezino ime pojavljuje u deklaraciji
prijateljske klase.
Ipak, važno je naglasiti da bi u ovom slučaju 
uključivanje datoteke zaglavlja \textit{Kredit.h}
u datoteci \textit{Klijent.h} bilo pogrešno
zbog pojave \textbf{kružne ovisnosti}
(engl.\ \textit{circular dependency}).
Preciznije,
dobili bismo situaciju da datoteka \textit{Klijent.h}
ovisi o datoteci \textit{Kredit.h}
(jer ju uključuje) te da
datoteka \textit{Kredit.h} ovisi o datoteci \textit{Klijent.h}.
Napomenimo i da ovdje ne pomaže upotreba direktive
\begin{minted}{C++}
#pragma once
\end{minted}
u svakoj od datoteka zaglavlja jer ona samo sprječava višestruko uključivanje istog zaglavlja
(ne rješava problem kad u obradi prve datoteke kompajler još nema potrebnu informaciju iz druge datoteke,
a ta druga pak datoteka zahtijeva obradu te prve datoteke).
Jedno moguće rješenje za pojavu kružne ovisnosti
(koje ćemo koristiti i u sljedećem zadatku)
je upotreba deklaracije klase unaprijed.
\medskip

\textbf{Deklaracija unaprijed} (engl.\ \textit{forward declaration}) 
govori kompajleru da određena klasa, struktura, funkcija ili tip postoji, bez prikazivanja cijele definicije.
Primjerice,
ako neka klasa \texttt{Profesor} zahtijeva postojanje klase \texttt{Kolegij},
bez poznavanja članova koje ta klasa ima,
mogli bismo u datoteci zaglavlja klase \texttt{Profesor} (bez potrebe za uključivanjem datoteke zaglavlja te klase \texttt{Kolegij})
naznačiti komapjleru da ta klasa \texttt{Kolegij} postoji na sljedeći način:
\begin{minted}{C++}
class Kolegij; // deklaracija unaprijed klase Kolegij
class Profesor {
    Kolegij* kolegiji;
    ...
};
\end{minted}
Međutim, to radi samo ukoliko koristimo u klasi \texttt{Profesor}
pokazivače ili reference na objekte tipa klase \texttt{Kolegij},
te ne koristimo članove klase \texttt{Kolegij}.
Primjerice, sljedeće bi bilo neispravno:
\begin{minted}{C++}
class Kolegij; // deklaracija unaprijed klase Kolegij
class Profesor {
friend int Kolegij::sifra_kolegija();
...
};
\end{minted}


\begin{zadatak}
    U zadatku \ref{zad:kredit} u klasi \texttt{Klijent}
    deklarirali smo klasu \texttt{Kredit} kao njoj prijateljsku klasu.
    Međutim, ne trebaju sve funkcije članice klase \texttt{Kredit}
    pristup privatnim članovima klase \texttt{Klijent}.
    Izmijenite potrebne dijelove odgovarajućih datoteka zaglavlja tih klasa
    kako bi klasa \texttt{Klijent} umjesto deklariranja
    prijateljske klase \texttt{Kredit},
    deklarirala samo one metode klase \texttt{Kredit}
    kao prijateljske koje su potrebne za ispravno kompajliranje koda.
\end{zadatak}


\begin{rjesenje}
    U klasi \texttt{Kredit} jedino funkcija članica
    \texttt{provjera\_sposobnosti} pristupa dijelovima klase \texttt{Klijent}
    koji nisu javni.
    Prema tome, možemo samo nju napraviti prijateljskom funkcijom klase \texttt{Klijent} (a ne cijelu klasu \texttt{Kredit}):
\begin{minted}[
        linenos,
        numbersep=5pt
    ]{C++}
class Klijent {
// friend class Kredit; - ovu liniju koda ukloniti
friend bool Kredit::provjera_sposobnosti(Klijent&);
...
};
\end{minted}
Kako koristimo članove klase \texttt{Kredit}
(preciznije, metodu \texttt{provjera\_sposobnosti}),
potrebno je u datoteku \textit{Klijent.h} dodati
\begin{minted}[
        linenos,
        numbersep=5pt
    ]{C++}
#include "Kredit.h"
\end{minted}
No, kako bismo sada izbjegli pojavu kružne ovisnosti
(obzirom da datoteka \textit{Kredit.h} uključuje upravo izmijenjenu
\textit{Klijent.h} datoteku),
u datoteci \textit{Klijent.h} zamijenimo
uključivanje datoteke \texttt{Kredit.h} sa deklaracijom unaprijed klase \texttt{Kredit}:
\begin{minted}[
        linenos,
        numbersep=5pt
    ]{C++}
//#include "Klijent.h" - ovu liniju koda ukloniti

class Klijent; // dodana deklaracija unaprijed

class Kredit {
    ...
};
\end{minted}
\end{rjesenje}



\section{Nadogradnja operatora}


Sadržaj \textit{Matrica.h} datoteke:

\begin{minted}{C++}
#pragma once
#include <iostream>
#include <utility>

class Matrica {
public:
	Matrica(size_t a = 0,size_t b = 0,double c = 0);
	~Matrica();
	Matrica(const Matrica &);
private:
	void alociraj();
	void dealociraj();
	std::pair<size_t,size_t> dim;
	double **elementi;
};
\end{minted}

Sadržaj \textit{Matrica.cpp} datoteke:

\begin{minted}{C++}
#include <iostream>
#include <utility>
#include "Matrica.h"

using namespace std;

void Matrica::alociraj() {
	elementi = new double*[dim.first];
	for(size_t i = 0; i < dim.first; ++i)
		elementi[i] = new double[dim.second];
}

void Matrica::dealociraj() {
	for(size_t i = 0; i < dim.first; ++i)
		delete[] elementi[i];
	delete[] elementi;
	elementi = nullptr;
}

Matrica::Matrica(size_t a, size_t b, double c) {
	dim.first = a;
	dim.second = b;
	alociraj();
	for(size_t i = 0; i < dim.first; ++i)
		for(size_t j = 0; j < dim.second; ++j)
			elementi[i][j] = c;
}

Matrica::~Matrica() {
	dealociraj();
}

Matrica::Matrica(const Matrica &m) {
	dim = m.dim;
	alociraj();
	for(size_t i = 0; i < dim.first; ++i)
		for(size_t j = 0; j < dim.second; ++j)
			elementi[i][j] = m.elementi[i][j];
}
\end{minted}

Sadržaj \textit{main.cpp} datoteke:

\begin{minted}{C++}
#include <iostream>
#include "Matrica.h"

using namespace std;

int main(void) {
	Matrica A(2,3,1.5), 
		B(2,3,2.4),
		C(2,3), D;
	
	return 0;
}
\end{minted}


\textbf{Nadogradnja} ili \textbf{preopterećivanje} operatora (engl. \textit{operator overloading}) 
omogućuje definiranje značenja postojećih C++ operatora kada se primjenjuju na objekte korisnički definiranih klasa. 
Na taj način objekti naših klasa mogu se koristiti slično ugrađenim tipovima podataka, što čini programe preglednijima i jednostavnijima za pisanje.
\medskip

Primjerice, za ugrađene numeričke tipove prirodno je pisati izraze poput:
\begin{minted}{C++}
cout << 2 + 3 * 4 << endl;
\end{minted}

Kada radimo s vlastitim klasama, primjerice klasom \texttt{Matrica}
koja reprezentira matrice realnih brojeva, željeli bismo zadržati isti način zapisivanja izraza:
\begin{minted}{C++}
cout << A + 2 * B << endl;
\end{minted}

Bez nadogradnje operatora morali bismo koristiti posebne funkcije poput:
\begin{minted}{C++}
ispis(zbroj(A, umnozak(2, B)));
\end{minted}

Iako bi oba zapisa dala isti rezultat, korištenje operatora znatno povećava čitljivost programa.
Važno je naglasiti da se nadogradnjom operatora ne stvaraju novi operatori niti se mijenja njihova sintaksa, prioritet ili broj operanada. 
Nadogradnjom samo definiramo kako će postojeći operatori djelovati nad objektima određene klase.
\medskip

Općenito, nadograđeni operator implementira se kao funkcija čije je ime oblika \texttt{operatorX}, 
pri čemu je \texttt{X} simbol operatora koji nadograđujemo.
Broj i tip operanada, kao i tip povratne vrijednosti,
ovisi o operatoru \texttt{X} koji nadograđujemo.
U nastavku ove točke proći ćemo operatore koji se najčešće nadograđuju.


\subsection{Nadogradnja operatora za unos i ispis}


Primjerice, za objekte naše vlastite klase \texttt{Matrica}
želimo pomoću sljedeće linije koda omogućiti ispis objekata \texttt{A} i \texttt{B}
tog tipa \texttt{Matrica}:

\begin{minted}{C++}
cout << A << B << endl;
\end{minted}

Primijećujemo da operator \texttt{<\!<} ima dva operanda,
jedan s lijeve strane i jedan s desne strane simbola \texttt{<\!<} - 
objekt \texttt{std::cout} (tipa \texttt{std::ostream})
i objekt tipa \texttt{Matrica}.
Povratna vrijednost operatora \texttt{<\!<} mora biti referenca na objekt s lijeve strane, 
tj.\ na \texttt{std::cout} kako bi ulančavanje
poziva operatora \texttt{<\!<} u gornjem primjeru
bilo moguće.
Preciznije, gornji izraz je jednak, gledajući prioritete,
izrazu 
\begin{minted}{C++}
((cout << A) << B) << endl;
\end{minted}
pa izraz u unutarnjoj zagradi treba vratiti referencu na 
objekt \texttt{cout} kako bi se na njega zatim
mogao ispisati objekt \texttt{B}.
\medskip

Slično, želimo pomoću operatora za unos \texttt{>\!>}
omogućiti unos objekata \texttt{A} i \texttt{B}
tipa naše vlastite klase \texttt{Matrica} kao u sljedećem primjeru:

\begin{minted}{C++}
cin >> A >> B;
\end{minted}

Operator \texttt{>\!>} ima dva operanda - objekt \texttt{std::cin} (tipa \texttt{std::istream})
i objekt tipa \texttt{Matrica}.
Ponovo, zbog ulančavanja je povratna vrijednost referenca na objekt s lijeve strane tog operatora.
\medskip

Istaknimo da se objekti tipa \texttt{std::ostream} i \texttt{std::istream} ne mogu kopirati
(pa prema tome koristimo reference u argumentima funkcija i njihovim povratnim vrijednostima
koje se odnose na te tipove).


\begin{zadatak}\label{zad:operator-unos-i-ispis}
Nadogradite operatore \texttt{>\!>} i \texttt{<\!<} za klasu
\texttt{Matrica} iz \textcolor{red}{zadatka ...} 
kako bi se objekti te klase mogli unositi
i ispisivati pomoću izraza \texttt{cin >\!> A} i \texttt{cout <\!< A}
(pri čemu je objekt \texttt{A} tipa klase \texttt{Matrica}.
Napišite i program za testiranje koji učitava elemente
dviju matrica te ih ispisuje upotrebom
nadograđenih operatora.
\end{zadatak}


\begin{rjesenje}
Matrice ćemo ispisati na način da ispišemo njihove elemente po retcima matrice.
Svaki redak matrice bit će ispisan u zasebnom retku, a
elementi pojedinog retka će odvojeni jednim znakom razmaka.
Matrice unosimo tako da unosimo element po element matrice, također po retcima.
\medskip

Dodamo u \textit{Matrica.cpp} datoteku implementaciju nadograđenog operatora za ispis matrica
te nadograđenog operatora za unos matrica:
\begin{minted}[
        linenos,
        numbersep=5pt
    ]{C++}
ostream& operator<<(ostream& os, const Matrica &m) {
    for(size_t i = 0; i < m.dim.first; ++i) {
        for(size_t j = 0; j < m.dim.second; ++j) {
            os << m.elementi[i][j] << " ";
        }
        cout << endl;
    }
    return os;
}

istream& operator>>(istream &is, Matrica &m) {
    for(size_t i = 0; i < m.dim.first; ++i) {
        for(size_t j = 0; j < m.dim.second; ++j) {
            is >> m.elementi[i][j];
        }
    }
    return is;
}
\end{minted}
Uočimo da u implementaciji operatora \texttt{>\!>}
matrica \texttt{m} nije označena sa \texttt{const}
obzirom da ju ta funkcija treba promijeniti.
Slično, unos i ispis podataka mijenjaju stanje odgovarajućih \texttt{ostream}
i \texttt{istream} objekata korištenih u tim funkcijama pa niti njih ne možemo označiti sa \texttt{const}.
\medskip

Obzirom da prikazane verzije nadograđenih operatora pristupaju
privatnim članovima klase \texttt{Matrica},
deklariramo ih kao prijateljske funkcije u toj klasi
(u \textit{Matrica.h} datoteci):
\begin{minted}[
        linenos,
        numbersep=5pt
    ]{C++}
class Matrica {
friend std::ostream& operator<<(std::ostream& os, 
    const Matrica& m);
friend std::istream& operator>>(std::istream& is, Matrica& m);
...
};
\end{minted}

Sadržaj datoteke \textit{main.cpp} s programom za testiranje:
\begin{minted}[
        linenos,
        numbersep=5pt
    ]{C++}
#include <iostream>
#include "Matrica.h"

using namespace std;

int main(void) {
    Matrica A(2,3), B(2,3);
    cout << "Unesite elemente matrice A: ";
    cin >> A;
    cout << "Unesite elemente matrice B: ";
    cin >> B;

    cout << "Matrica A: " << A << endl
         << "Matrica B: " << B << endl;

    return 0;
}
\end{minted}
\end{rjesenje}

Operatori su posebne funkcije, pa se kao i ostale
funkcije mogu izravno pozivati.
Tako bismo, primjerice, mogli za objekte \texttt{A} i \texttt{B} tipa klase \texttt{Matrica} umjesto
\begin{minted}{C++}
cout << A << B << endl;
\end{minted}
pisati
\begin{minted}{C++}
operator<<(operator<<(cout,A), B) << endl;
\end{minted}
Namjerno smo u gornjem izdvojili \texttt{endl}
iz izravnog poziva funkcije \texttt{operator<\!<}
kako bi naglasili da se tu ne primjenjuje naša
verzija nadograđenog operatora \texttt{<\!<}
(obzirom da se \texttt{endl} ne može pretvoriti u tip klase \texttt{Matrica}).


\subsection{Nadogradnja aritmetičkih operatora}

Operatori za zbrajanje, oduzimanje, množenje, dijeljenje i modulo (\texttt{+}, \texttt{-}, \texttt{*}, \texttt{/} i \texttt{\%}).

\begin{zadatak}\label{zad:matrice-aritmetika}
    Nadogradite operator \texttt{+} za zbrajanje objekata klase \texttt{Matrica}
    iz zadatka \ref{zad:operator-unos-i-ispis}.
    Nadogradite i operator \texttt{*}
    tako da izraz \texttt{2 * A} predstavlja množenje matrice skalarom (pri čemu je skalar podatak tipa \texttt{double}).
    Sljedeći kod bi se trebao uspješno kompajlirati i izvršavati:
    \begin{minted}[
        linenos,
        numbersep=5pt
    ]{C++}
        Matrica A(2,3), B(2,3);
        cin >> A >> B;
        cout << A + B << endl << 2 * A;
    \end{minted}
\end{zadatak}


\begin{rjesenje}
    Operator za zbrajanje kao argumente ima dvije matrice koje zbrajamo
    (koristimo reference kako ne bi nepotrebno kopirali te matrice,
    te koristimo \texttt{const} obzirom da operator ne mijenja niti jednu od tih matrica).
    Rezultat je nova matrica koju funkcija vraća kao rezultat.
    Preciznije, zbroj dviju matrica izračunat ćemo u lokalnoj varijabli (tipa klase \texttt{Matrica}) te ćemo vratiti kopiju te lokalne varijable.
    \medskip

    Operator množenja \texttt{*} u ovom slučaju koristimo tako da s lijeve strane simbola \texttt{*} napišemo realan broj, a 
    s desne strane matricu.
    Tako da će operator u ovom slučaju imati kao prvi argument realan broj tipa \texttt{double}, a kao drugi element matricu.
    Operator ne mijenja niti realan broj niti matricu,
    a rezultat je nova matrica koju će funkcija vratiti.
    \medskip

    Prema tome, u datoteku \texttt{Matrica.cpp} dodajemo 
    potrebne implementacije:
    \begin{minted}[
        linenos,
        numbersep=5pt
    ]{C++}
Matrica operator+(const Matrica &lm, const Matrica &dm) {
    Matrica zbroj(lm); // iskoristimo konstruktor kopiranjem
    for(size_t i = 0; i < zbroj.dim.first; ++i) {
        for(size_t j = 0; j < zbroj.dim.second; ++j) {
            zbroj.elementi[i][j] += dm.elementi[i][j];
        }
    }
    return zbroj;
}

Matrica operator*(const double &lbr, const Matrica &dm) {
    Matrica rezultat(dm);
    for(size_t i = 0; i < dm.dim.first; ++i)
        for(size_t j = 0; j < dm.dim.second; ++j)
            rezultat.elementi[i][j] *= lbr;
    return rezultat;
}
    \end{minted}

Obzirom da navedene funkcije koriste privatne članove klase \texttt{Matrica},
dodamo u tu klasu (u \textit{Matrica.h} datoteci)
njihove deklaracije kao deklaracije prijateljskih funkcija te klase:

\begin{minted}[
        linenos,
        numbersep=5pt
    ]{C++}
class Matrica {
friend Matrica operator+(const Matrica& lm, 
    const Matrica& dm);
friend Matrica operator*(const double& lbr, 
    const Matrica& dm);
...
};
\end{minted}
\end{rjesenje}


\begin{razmisljanje}
    \begin{enumerate}
        \item 
        Zašto nije ispravno u, primjerice, našoj nadogradnji operatora \texttt{+}
        vratiti referencu, tj.\ zašto sljedeći kod nije ispravan:

        \begin{minted}{C++}
Matrica& operator+(const Matrica &lm, const Matrica &dm) 
{
    Matrica zbroj(lm);
    for(size_t i = 0; i < zbroj.dim.first; ++i) {
        for(size_t j = 0; j < zbroj.dim.second; ++j) {
            zbroj.elementi[i][j] += dm.elementi[i][j];
        }
    }
    return zbroj;
}
    \end{minted}

    \item 
    Što biste promijenili u našoj nadogradnji operatora \texttt{+} iz prethodnog zadatka
    kako bi dobili nadogradnju operatora \texttt{-} za oduzimanje matrica?
    Kako biste implementirali operator \texttt{*} za množenje dvije matrice?

    \item 
    Kako biste, za dvije matrice \texttt{A} i \texttt{B},
    nadogradili operator \texttt{-} oduzimanja matrica
    upotrebom operatora \texttt{+} za zbrajanje
    matrica i operatora \texttt{*} za množenje matrice skalarom?

    \item 
    Kako biste napisali unarni operator \texttt{-} za klasu \texttt{Matrica}?
    Koliko argumenata takav operator ima i koja mu je povratna vrijednost?

    \item 
    Iako mi to radi jednostavnosti koda ovdje nismo radili,
    u nekim prethodno spomenutim operatorima nužne
    su provjere matrica koje su im proslijeđene kao argumenti.
    Primjerice, dvije matrice možemo zbrojiti samo ako su one istih dimenzija.
    Koje biste provjere dodali u ostale spomenute operatore?
    Kako biste prekinuli izvršavanje programa u slučaju da, primjerice,
    operatoru za zbrajanje matrica budu kao argumenti proslijeđene matrice različitih dimenzija?
    \end{enumerate}
\end{razmisljanje}


\subsection{Nadogradnja relacijskih operatora}


Relacijski operatori \texttt{==}, \texttt{!=},
\texttt{<}, \texttt{<=}, \texttt{>}, \texttt{>=}
najčešće kao povratnu vrijednost imaju podatak tipa \texttt{bool}
(jesu li argumenti u navedenoj relaciji ili nisu).
Operator \texttt{==} za testiranje jednakosti dva objekta neke klase 
obično se svodi na
provjeru jesu li odgovarajući dijelovi objekata tipa te klase jednaki.
Također, korisnik u slučaju da je nadograđen operator \texttt{==}
obično očekuje da je nadograđen i operator \texttt{!=}
(pri čemu se on definira pomoću tog, već nadograđenog, operatora \texttt{==}).
Slično se često očekuje da definicije operatora 
\texttt{<}, \texttt{<=}, \texttt{>} i \texttt{>=}
budu konzistentne međusobno i s definicijama operatora \texttt{==} i \texttt{!=}.
Primjerice, ako su dva objekta različita u smislu operatora \texttt{!=}
tada mora biti jedan manji (\texttt{<}) ili veći (\texttt{>}) od drugoga
(to je karakteristično za linearno ili totalno uređene skupove,
poput skupa realnih brojeva).


\begin{zadatak}\label{zad:matrice-usporedba}
    Za klasu \texttt{Matrica} iz zadatka \ref{zad:matrice-aritmetika}
    nadogradite operator usporedbe \texttt{==}
    te operator \texttt{!=}.
    Za matrice \texttt{A} i \texttt{B},
    izraz \texttt{A == B} jednak je \texttt{true} ako i samo ako su matrice \texttt{A} i \texttt{B} jednake,
    dok je izraz \texttt{A == B} jednak \texttt{true} ako i samo ako su matrice \texttt{A} i \texttt{B} različite.

     Sljedeći kod bi se trebao uspješno kompajlirati i izvršavati:
    \begin{minted}[
        linenos,
        numbersep=5pt
    ]{C++}
        Matrica A(2,3), B(2,3);
        cin >> A >> B;
        cout << (A == B) << endl << (A != 2 * B) << endl;
    \end{minted}
\end{zadatak}


\begin{rjesenje}
    Dvije matrice biti će jednake ako i samo ako imaju jednake
    dimenzije i jednake odgovarajuće elemente.
    Operator \texttt{!=} jednostavno ćemo implementirati upotrebom operatora \texttt{==}
    kojeg smo prije njega nadogradili
    upotrebom činjenice da su dvije matrice jednake
    ako i samo ako ne vrijedi da su jednake.
    \medskip

    Dodajemo nadogradnje tih operatora u \textit{Matrica.cpp} datoteku:
    \begin{minted}[
        linenos,
        numbersep=5pt
    ]{C++}
bool operator==(const Matrica &lm, const Matrica &dm) {
    if(lm.dim != dm.dim) // uspoređivanje parova
        return false;
    for(size_t i = 0; i < lm.dim.first; ++i) {
        for(size_t j = 0; j < lm.dim.second; ++j) {
            if(lm.elementi[i][j] != dm.elementi[i][j]) {
                return false;
            }
        }
    }
    return true;
}

bool operator!=(const Matrica &lm, const Matrica &dm) {
    return !(lm == dm);
}
    \end{minted}

Kako te funkcije koriste privatne članove klase \texttt{Matrica},
deklariramo ih u toj klasi (u \textit{Matrica.h} datoteci) kao prijateljske funkcije te klase:
\begin{minted}[
        linenos,
        numbersep=5pt
    ]{C++}
class Matrica {
friend bool operator==(const Matrica &lm, const Matrica &dm);
friend bool operator!=(const Matrica &lm, const Matrica &dm);
...
};
\end{minted}
\end{rjesenje}


\begin{razmisljanje}
    \begin{enumerate}
        \item
        Kako biste definirali neki uređaj \texttt{<} na matricama?
        \item 
        Ako za matrice (klase \texttt{Matrica})
        nadogradimo operator \texttt{<},
        kako biste pomoću njega (i operatora koje smo za tu klasu prethodno nadogradili)
        nadogradili preostale relacijske operatore
        (\texttt{<=}, \texttt{>}, \texttt{>=})?
    \end{enumerate}
\end{razmisljanje}


\subsection{Nadogradnja operatora pridruživanja}

Operator pridruživanja \texttt{=} implementiramo kao funkciju članicu
klase na koju se odnosi.
Kod takve implementacije, prvi operand (onaj s lijeve strane
operatora pridruživanja \texttt{=}) vezan je uz implicitni pokazivač \texttt{this}.
To znači da funkcija koju pišemo ima jedan operand manje (preciznije,
imamo samo jedan operand i to onaj koji se nalazi s desne strane operatora pridruživanja \texttt{=}).
Kao povratni tip vraćamo referencu na objekt s lijeve strane (onaj na koji se odnosi implicitan pokazivač \texttt{this})
jer želimo omogućiti ulančavanje,
poput \texttt{A = B = C}
(ovdje se prvo pridružuje \texttt{B = C} i rezultat tog izraza je referenca na objekt \texttt{B}
koji se zatim pridružuje objektu \texttt{A}).


\begin{zadatak}\label{zad:matrice-pridruzivanje}
Nadogradite operator pridruživanja za klasu \texttt{Matrica} iz zadatka \ref{zad:matrice-usporedba}.
Pripazite da operator ispravno radi u slučaju samopridruživanja,
tj.\ kad za neku matricu \texttt{A} imamo izraz \texttt{A = A}).
\medskip

Sljedeći kod bi se trebao uspješno kompajlirati i izvršavati:
    \begin{minted}[
        linenos,
        numbersep=5pt
    ]{C++}
        Matrica A(2,3), B(5,6);
        cin >> A;
        A = A = B;
        cout << B;
    \end{minted}
\end{zadatak}


\begin{rjesenje}
    Obzirom da će nadograđeni operator pridruživanja \texttt{=} biti funkcija članica klase \texttt{Matrica},
    dodajemo njegovu deklaraciju u tu klasu
    (u datoteci \textit{Matrica.h}):
    \begin{minted}[
        linenos,
        numbersep=5pt
    ]{C++}
class Matrica {
...
public:
    Matrica &operator=(const Matrica& dm);
...
};
    \end{minted}

    U datoteku \textit{Matrica.cpp} dodajemo implementaciju te funkcije članica klase \texttt{Matrica}.
    Pritom prvo provjerimo za slučaj \texttt{A = A},
    je li objekt s lijeve strane (spremljen na adresi \texttt{this})
    isti kao i objekt
    na desnoj strani znaka \texttt{=}
    (prema tome na istoj adresi kao i \texttt{this}).
    Ako je to slučaj (tj.\ \texttt{A = A}),
    operator pridruživanja \texttt{=} zapravo ne treba raditi ništa.

    \begin{minted}[
        linenos,
        numbersep=5pt
    ]{C++}    
Matrica &Matrica::operator=(const Matrica& dm) {
    if(this != &dm) {
        dealociraj();
        dim = dm.dim;
        alociraj();
        for(size_t i = 0; i < dim.first; ++i) {
            for(size_t j = 0; j < dim.second; ++j) {
                elementi[i][j] = dm.elementi[i][j];
            }
        }
    }
    return *this;
}
    \end{minted}
\end{rjesenje}


\begin{razmisljanje}
    \begin{enumerate}
        \item 
         Ako za klasu \texttt{Matrica}
         iz prethodnog zadatka ne nadogradimo operator pridruživanja \texttt{=}
         tada će se sljedeći program kompajlirati, ali pri pokretanju srušiti:
        \begin{minted}{C++}
int main(void) {
    Matrica A(2,3), B(2,3);
    cin >> A;
    B = A;
    cout << B;

    return 0;
}
        \end{minted}
    Obrazložite.

    \item 
    Ako za klasu \texttt{Matrica}
         iz prethodnog zadatka ne nadogradimo operator pridruživanja \texttt{=},
         obrazložite zašto sljedeći program radi ispravno:
    \begin{minted}{C++}
int main(void) {
    Matrica A(2,3);
    cin >> A;
    Matrica B = A;
    cout << B;

    return 0;
}
        \end{minted}
    \end{enumerate}
\end{razmisljanje}


\subsection{Nadogradnja operatora složenog pridruživanja}

Operatori složenog pridruživanja (\texttt{+=}, \texttt{-=}, \texttt{*=}, \texttt{/=}, \texttt{\%=}) 
mijenjaju objekt s lijeve strane operatora
te je, kao i za operator pridruživanja \texttt{=}, 
uobičajeno da vraćaju referencu na taj objekt.
Zbog toga su najčešće implementirani kao funkcije članice klase objekta s lijeve strane operatora,
iako se od njih to ne zahtijeva
(mogu biti implementirani kao nečlanske funkcije,
pri čemu je tada često u odgovarajućoj klasi 
potrebna njihova deklaracija kao prijateljska funkcija toj klasi).


\begin{zadatak}\label{zad:matrice-slozeno-pridruzivanje}
    Nadogradite operator složenog pridruživanja \texttt{+=} za klasu \texttt{Matrica} iz zadatka \ref{zad:matrice-pridruzivanje}.
    Operator matricu s desne strane operatora \texttt{+=} pribraja matrici s lijeve strane tog operatora.
    \medskip

    Sljedeći kod bi se trebao uspješno kompajlirati i izvršavati:
    \begin{minted}[
        linenos,
        numbersep=5pt
    ]{C++}
        Matrica A(2,3), B(2,3);
        cin >> A >> B;
        (A += B) += A;
        cout << A;
    \end{minted}
\end{zadatak}


\begin{rjesenje}
    Obzirom da će nadograđeni operator \texttt{+=} biti funkcija članica klase \texttt{Matrica},
    dodajemo njegovu deklaraciju u tu klasu
    (u datoteci \textit{Matrica.h}):
    \begin{minted}[
        linenos,
        numbersep=5pt
    ]{C++}
class Matrica {
...
public:
    Matrica &operator+=(const Matrica& dm);
...
};
    \end{minted}
    U datoteku \textit{Matrica.cpp} dodajemo implementaciju te funkcije članica klase \texttt{Matrica}.
    Elementima matrice čija je to funkcija članica (tj.\ 
    matrici s lijeve strane operatora \texttt{+=}),
    dodajemo odgovarajuće elemente matrice s desne strane tog operatora
    (tj.\ matrice koja je primljena kao argument funkcije).
    Povratna vrijednost je referenca na promijenjenu matricu s lijeve strane operatora:
    \begin{minted}[
        linenos,
        numbersep=5pt
    ]{C++}
Matrica& Matrica::operator+=(const Matrica& dm) {
    for(size_t i = 0; i < dim.first; ++i)
        for(size_t j = 0; j < dim.second; ++j)
            elementi[i][j] += dm.elementi[i][j];
    return *this;
}
    \end{minted}
\end{rjesenje}



\begin{razmisljanje}
    \begin{enumerate}
        \item 
        U prethodnom zadatku u testnom primjeru navedena je sljedeća linija koda:
\begin{minted}{C++}
(A += B) += A;
\end{minted}
    Obrazložite razliku u posljedicama izvrašavanja te linije u odnosu na sljedeću liniju koda:
\begin{minted}{C++}
A += B += A;
\end{minted}

    \item 
    Prototip nadograđenog operatora iz rješenja prethodnog zadatka je
    \begin{center}
        \texttt{Matrica\& Matrica::operator+=(const Matrica\& dm)}
    \end{center}
    Obrazložite kako je uz ovakav prototip moguće kompajlirati i izvršiti sljedeću liniju koda
    \begin{minted}{C++}
    A += A;
    \end{minted}
    (ako prema toj liniji mijenjamo matricu s desne strane operatora koja je označena u nadograđenom operatoru kao \texttt{const}).

    \item 
    Zbog jednostavnosti u rješenju prethodnog zadatka nismo promatrali slučaj kad se s lijeve i desne strane operatora \texttt{+=}
    nalaze matrice različitih dimenzija.
    Kako biste u tom slučaju prekinuli izvrašavanje programa?
    Kako biste implementirali operatore složenog pridruživanja \texttt{-=} i \texttt{*=}
    za matrice?

    \item 
    Kako biste nadogradili operator složenog pridruživanja \texttt{*=} tako da 
    omogućuje množenje matrice skalarom
    (realan broj tipa \texttt{double})
    pomoću izraza \texttt{A *= 1.5}?
    \end{enumerate}
\end{razmisljanje}



\subsection{Nadogradnja operatora indeksiranja}


Klase koje predstavljaju spremnike podataka, odnosno strukture podataka iz kojih se elementi dohvaćaju prema njihovoj poziciji, 
najčešće definiraju operator indeksiranja \texttt{[]}. 
Njegova je svrha omogućiti pristup elementima na način sličan radu s poljima.
\medskip

Operator indeksiranja mora biti implementiran kao funkcija članica klase, jer se uvijek primjenjuje nad konkretnim objektom spremnika. 
Uobičajena implementacija vraća referencu na dohvaćeni element, što omogućuje ne samo čitanje njegove vrijednosti nego i njezino izravno mijenjanje.
Radi ispravne podrške za konstantne objekte, preporučuje se definirati i \texttt{const} verziju tog operatora,
tj.\ verziju koja se primjenjuje na konstantne objekte te vraća referencu na konstantan element
(čime se onemogućava njegova promjena).
\medskip

Primjerice,
ako klasa \texttt{Niz} sprema niz realnih brojeva tipa \texttt{double},
želimo nadogradnjom operatora indeksiranja \texttt{[]}
omogućiti pristup i promjenu pojedinog realnog broja u tom nizu na sljedeći način:
\begin{minted}{C++}
Niz nekonstantan;
nekonstantan[4] = 10.5;

const Niz konstantan;
cout << konstantan[1] << endl;
\end{minted}

Pritom se u liniji gornjeg koda pod brojem 2 koristi nekonstantna verzija operatora indeksiranja
(obzirom da želimo promijeniti element tog niza na \texttt{10.5}),
dok se u liniji pod brojem 5 koristi konstantna (tj.\ \texttt{const}) verzija
tog operatora
(jer se objekt \texttt{konstantan}
ne smije mijenjati obzirom da je označen kao \texttt{const}).

\begin{zadatak}
    ...
\end{zadatak}